% !PRINTED 301
\weberpage{301}

Es mögen also für einen Werth $\xi$ von $x$ zwischen $\alpha$ und $\beta$ in der Reihe

\weq{3}{f^{(\nu)}(x),\ f^{(\nu+1)}(x),\ .\ .\ .\ f^{(\nu+\mu-1)}(x),\ f^{(\nu+\mu)}(x)}

alle Glieder, mit Ausnahme des letzten, verschwinden, und das letzte $f^{(\nu+\mu)}(x)$ möge etwa einen positiven Werth haben.

Wir grenzen um $\xi$ zwei Intervalle $\delta_1$, $\delta_2$ ab, so dass alle Werthe von $x$ im Intervall $\delta_1$ kleiner, im Intervall $\delta_2$ grösser als $\xi$ sind, und nehmen diese Intervalle so klein, dass die Functionen (3) ausser in $\xi$ darin nicht verschwinden, also auch $f^{(\nu+\mu)}(x)$ positiv bleibt.

Da nun, wenn $f^{(\nu+\mu)}(x)$ positiv ist, $f^{(\nu+\mu-1)}(x)$ mit $x$ zugleich wächst, so ist

\[
\begin{array}{l}
f^{(\nu+\mu-1)}(x) \text{ in } \delta_1 \text{ negativ,} \\
\text{in } \delta_2 \text{ positiv.}
\end{array}
\]

Daraus folgt, dass $f^{(\nu+\mu-2)}(x)$ in $\delta_1$ abnimmt, in $\delta_2$ wächst, also in beiden Intervallen positiv ist, und so schliessen wir weiter auf die Vorzeichenfolge in der Reihe (3):

\[
\begin{array}{l}
(4)\quad f^{(\nu)}(x),\ f^{(\nu+1)}(x),\ f^{(\nu+2)}(x),\ .\ .\ .\ f^{(\nu+\mu-1)}(x),\ f^{\nu+\mu}(x) \\[4pt]
\delta_1,\quad (-1)^\mu \quad -(-1)^\mu \quad (-1)^\mu \quad .\ .\ . \quad - \quad + \\[2pt]
\delta_2,\quad + \quad + \quad + \quad .\ .\ . \quad + \quad +
\end{array}
\] % [sic] last term of row (4) prints as $f^{\nu+\mu}(x)$, without parentheses around the exponent, unlike the other terms in the same row

d.\,h. in der Reihe (3) werden beim Durchgang durch $\xi$ aus lauter Zeichenwechseln Zeichenfolgen, und es gehen in der Reihe (3) $\mu$ Zeichenwechsel verloren.

Ist $f^{(\nu+\mu)}(\xi)$ negativ, so sind alle Zeichen in (4) die entgegengesetzten, und der Schluss bleibt sonst derselbe.

Wenn nun $\nu = 0$, d.\,h. $f^{(\nu)}(x)$ die ursprüngliche Function $f(x)$ selbst ist, die dann in $\xi$ eine $(\mu+1)$fache Wurzel hat, so findet also beim Durchgang durch $\xi$ auch in der Reihe (1) ein Verlust von $\mu$ Zeichenwechseln statt.

Ist aber $\nu > 0$ und die dem $f^{(\nu)}(x)$ vorangehende Function $f^{(\nu-1)}(x)$ von Null verschieden, so haben wir folgende Zeichen:

1. $\mu$ gerade:

\[
\begin{array}{ll}
\text{a)}\ f^{(\nu-1)}(x),\ f^{(\nu)}(x), & \text{b)}\ f^{(\nu-1)}(x),\ f^{(\nu)}(x) \\[2pt]
\delta_1,\quad + \quad + & \quad - \quad + \\[2pt]
\delta_2.\quad + \quad + & \quad - \quad +
\end{array}
\]
